\begin{theorem}[Teorema de las diferencias de Van der Corput]
    \label{teo:diferencias_weyl}
    Sea $(x_n)_{n\geq1}$ una sucesión de números reales. Si para cada $h=1,2,\dots,$
    la sucesión $(x_{n+h}-x_n)$ es equidistribuida módulo $1$, entonces la sucesión
    $(x_n)$ es equidistribuida módulo $1$.
\end{theorem}

\begin{proof}
    Sea $k\in \mathbb{Z}\setminus\{0\}$ y sea $f(n)=kx_n$, usando la estimación 
    (\ref{lem:vandercorput}), y teniendo en cuenta que para $h=0$ se tiene $e^{2\pi i 0}=1$.
    Obtenemos
    \[
        \left|\frac{1}{N}\sum_{n=1}^N e^{2\pi i k x_n}\right|^2 \leq
        \frac{4}{H} + \frac{4}{H} \sum_{h=1}^{H-1} \frac{N-h}{N} \left|\frac{1}{N-h}\sum_{n=1}^{N-h} e^{2\pi i k (x_{n+h}-x_n)} \right|
    \]

    Dado que para cada $h=1,2,\dots,$ la sucesión $(x_{n+h}-x_n)$ es equidistribuida módulo $1$,
    por el criterio de Weyl se tiene que
    \[
        \lim_{N\to\infty} \frac{1}{N-h}\sum_{n=1}^{N-h} e^{2\pi i k (x_{n+h}-x_n)} = 0.
    \]
    Luego,
    \[
        \limsup_{N\to\infty} \left|\frac{1}{N}\sum_{n=1}^N e^{2\pi i k x_n}\right|^2 \leq \frac{4}{H}.
    \]
    
    Dado que es cierto para cada $H\in \mathbb{N}$, haciendo $H\to\infty$ se concluye que
    \[
        \lim_{N\to\infty} \left|\frac{1}{N}\sum_{n=1}^N e^{2\pi i k x_n}\right|^2 = 0.
    \]

    Por lo tanto, por el criterio de Weyl, la sucesión $(x_n)$ es equidistribuida módulo $1$.
\end{proof}
    